\section*{Chapitre \ref{ch:b2:cell-signalling} --- Messagers chimiques et transduction du signal}

\begin{solution}{exo:b2:cell-signalling:1}
Insuline~: endocrine, hydrosoluble (peptide). Cortisol~: endocrine,
liposoluble. Acétylcholine à une synapse~: synaptique, hydrosoluble.
Histamine dans une plaie~: paracrine, hydrosoluble. Œstradiol~:
endocrine, liposoluble. Adrénaline de la surrénale~: endocrine,
hydrosoluble. Monoxyde d'azote~: paracrine, liposoluble (un gaz qui
traverse les membranes).
\end{solution}

\begin{solution}{exo:b2:cell-signalling:2}
L'adrénaline se fixe sur le récepteur (arrêt~: dissociation,
internalisation du récepteur)~; le récepteur active la \omterm{prop:b2:cell-signalling:gpcr}{protéine G}, qui
échange son GDP contre du GTP (arrêt~: hydrolyse du GTP en quelques
secondes)~; la \omterm{prop:b2:cell-signalling:gpcr}{protéine G} active l'adénylate cyclase, qui fabrique
l'AMPc (arrêt~: phosphodiestérase)~; l'AMPc active la kinase A (arrêt~:
disparition de l'AMPc)~; la kinase A phosphoryle la phosphorylase
kinase, qui phosphoryle la glycogène phosphorylase (arrêt~:
phosphatases)~; la phosphorylase libère du glucose 1-phosphate à partir
du glycogène, qui est converti et exporté sous forme de glucose.
\end{solution}

\begin{solution}{exo:b2:cell-signalling:3}
Les récepteurs de surface fixent les messagers hydrosolubles à la
membrane et agissent en quelques secondes par des \omterm{prop:b2:cell-signalling:gpcr}{seconds messagers} et
des kinases, en modifiant l'activité de protéines existantes. Les
\omterm{prop:b2:cell-signalling:nuclear}{récepteurs nucléaires} fixent les messagers liposolubles à l'intérieur de
la cellule et agissent en quelques heures comme \omterm{prop:b2:cell-differentiation:transcriptional}{facteurs de
transcription}, en modifiant les protéines fabriquées.
\end{solution}

\begin{solution}{exo:b2:cell-signalling:4}
\omterm{prop:b2:cell-signalling:gpcr}{AMP cyclique}~: fabriqué par l'adénylate cyclase, détruit par la
phosphodiestérase. Calcium~: entre par des canaux ou sort du réticulum
endoplasmique en réponse à l'inositol trisphosphate, éliminé par des
pompes. Inositol trisphosphate (avec le diacylglycérol)~: fabriqué par
la phospholipase C à partir d'un lipide membranaire, éliminé par des
phosphatases.
\end{solution}

\begin{solution}{exo:b2:cell-signalling:5}
$\theta = [L]/(K_d + [L])$~: $0.09$, $0.5$, $0.91$, $0.99$. Doubler le
nombre de récepteurs laisse la fraction inchangée et double le nombre de
récepteurs occupés --- et donc la réponse.
\end{solution}

\begin{solution}{exo:b2:cell-signalling:6}
$50\times 200\times 100\times 300 = 3\times 10^{8}$. Cent récepteurs
occupés donnent $3\times 10^{10}$ molécules de produit par seconde ---
une borne supérieure, puisque les étapes saturent.
\end{solution}

\begin{solution}{exo:b2:cell-signalling:7}
La \omterm{prop:b2:cell-signalling:gpcr}{protéine G} bloquée maintient la cyclase active, l'AMPc reste élevé,
la kinase A maintient phosphorylé et ouvert le canal chlorure des
cellules intestinales, le chlorure est sécrété, et le sodium et l'eau le
suivent dans la lumière --- des litres par jour. Le blocage est une
modification covalente de la \omterm{prop:b2:cell-signalling:gpcr}{protéine G}, qui ne disparaît que lorsque
les cellules elles-mêmes sont remplacées, des jours plus tard.
\end{solution}

\begin{solution}{exo:b2:cell-signalling:8}
$0.9\times 10^{-6}\,\text{mol/L}\times 2\times 10^{-12}\,\text{L} =
1.8\times 10^{-18}$ mol, soit $1.1\times 10^{6}$ ions (davantage en
réalité, puisque des tampons fixent l'essentiel de ce qui entre). Un
seul canal à $10^{6}$ ions par seconde mettrait une seconde~; mille
canaux le font en une milliseconde, l'échelle de temps dont un signal a
besoin.
\end{solution}

\begin{solution}{exo:b2:cell-signalling:9}
Les deux voies convergent sur les mêmes enzymes~: la kinase A du glucagon
phosphoryle la glycogène synthase (inactivée) et la phosphorylase kinase
(activée)~; la cascade de l'insuline active la phosphatase qui retire les
phosphates. L'état des enzymes est fixé par l'équilibre entre activité
kinase et activité phosphatase, si bien que c'est le rapport des deux
\omterm{def:b2:cell-signalling:kinds}{hormones} qui compte~; lorsque les deux sont élevées, la phosphatase
l'emporte en général et le glycogène est stocké.
\end{solution}

\begin{solution}{exo:b2:cell-signalling:10}
Adrénaline~: l'AMPc diffuse à travers une cellule de \qty{20}{\micro m}
en moins d'une seconde ($x^{2}/2D$ avec $D \approx \qty{3e-10}{m^2/s}$),
et toutes les enzymes de la cascade existent déjà --- quelques secondes.
Cortisol~: un transcrit demande \qty{60}{s}, une protéine \qty{100}{s}~;
à partir de quelques dizaines d'ARNm portant chacun une douzaine de
ribosomes, la cellule fabrique quelques centaines de molécules d'enzyme
par minute, si bien que les milliers nécessaires à un effet mesurable
demandent des heures, après les minutes qu'il faut au cortisol pour
entrer et atteindre l'ADN.
\end{solution}

\begin{solution}{exo:b2:cell-signalling:11}
Le récepteur signale en permanence~: la cascade des MAP kinases pousse la
cellule à se diviser sans aucun facteur de croissance, et la cellule
ignore l'absence du signal qui la retient normalement. Comme une seule
\omterm{def:b2:genome-diversification:mutation}{mutation} supprime un contrôle de la division, ces récepteurs (et les
kinases situées en aval) comptent parmi les oncogènes les plus
fréquents. Un médicament qui occupe le site de l'ATP de la kinase arrête
les phosphorylations et fait taire le récepteur --- le principe de
plusieurs médicaments anticancéreux.
\end{solution}

\begin{solution}{exo:b2:cell-signalling:12}
Endocrine~: atteint toutes les cellules en une minute, peu coûteux par
cellule, incapable de s'adresser à une seule cellule, agit pendant des
minutes à des jours --- indispensable aux états du corps entier (jeûne,
croissance, reproduction). Nerveux~: une milliseconde, une seule cible,
un câblage coûteux, agit pendant des millisecondes --- indispensable au
mouvement et aux réflexes rapides. Les deux servent au même message
lorsque la vitesse et la portée sont toutes deux requises~: lors d'une
frayeur, les nerfs sympathiques et la médullosurrénale délivrent tous
deux de la noradrénaline ou de l'adrénaline.
\end{solution}

\begin{solution}{pb:b2:cell-signalling:1}
\textbf{1.} $1\,\text{nmol}/5\,\text{L} = \qty{0.2}{nmol/L}$.
\textbf{2.} $\theta = 0.2/1.2 = 0.17$~: environ 3300 récepteurs occupés.
\textbf{3.} Trois demi-vies~: \qty{0.025}{nmol/L}~; $\theta = 0.024$.
\textbf{4.} Entre une demi-minute et une minute, le temps que met le
\omterm{def:b2:blood-circulation:blood}{sang} pour faire le tour du corps~; le nerf délivre son \omterm{def:b2:cell-signalling:kinds}{neurotransmetteur}
directement au foie en moins d'une seconde.
\textbf{5.} La même fraction, mais $33\,000$ récepteurs occupés~: une
réponse dix fois plus grande.
\textbf{6.} À \qty{1}{nmol/L}, un récepteur de $K_d = \qty{1}{\micro
mol/L}$ est occupé au millième~: aucune réponse. Accorder $K_d$ à la
gamme circulante place la partie raide de la courbe d'occupation là où
l'\omterm{def:b2:cell-signalling:kinds}{hormone} varie réellement.
\textbf{7.} $20\times 100\times 1\times 50\times 50\times 1000 =
5\times 10^{9}$ molécules de glucose par seconde par récepteur occupé.
\textbf{8.} $3300\times 20\times 100 = 6.6\times 10^{6}$ AMPc pendant
la première seconde~; dans \qty{5e-12}{L}, $1.1\times 10^{-17}$ mol,
soit environ \qty{2}{\micro mol/L}.
\textbf{9.} $3300\times 5\times 10^{9} = 1.7\times 10^{13}$ par seconde, $2\times
10^{13}$ par minute par cellule.
\textbf{10.} $2\times 10^{11}$ cellules $\times 10^{15} = 2\times
10^{26}$ molécules par minute, 330 mol, \qty{6e4}{g} par minute ---
dix mille fois les \qty{5}{g/min} observés. Les gains ne se multiplient
que tant qu'aucune étape n'est saturée~; en réalité, la phosphorylase est
limitée par son substrat et par les phosphatases, et l'exportation du
glucose par les transporteurs. Le gain de la cascade sert à la rapidité
et à la sensibilité, pas au débit.
\textbf{11.} À \qty{5}{g/min}, vingt minutes~; la frayeur est passée en
quelques minutes et les interrupteurs d'arrêt stoppent la cascade bien
avant.
\textbf{12.} L'hydrolyse du GTP par la \omterm{prop:b2:cell-signalling:gpcr}{protéine G} (secondes)~; la
phosphodiestérase sur l'AMPc (secondes)~; les phosphatases sur les
enzymes phosphorylées (secondes à minutes)~; la \omterm{prop:b2:reproductive-hormones:pulses}{désensibilisation des
récepteurs} (minutes). Les plus rapides sont l'horloge propre de la
\omterm{prop:b2:cell-signalling:gpcr}{protéine G} et la phosphodiestérase.
\textbf{13.} L'AMPc n'est pas détruit~: il s'accumule et persiste, la
kinase reste active, et la production de glucose est plus forte et dure
plus longtemps --- l'effet éveillant du café.
\textbf{14.} $4000/50 = \qty{80}{s}$.
\textbf{15.} $600/5 = \qty{120}{s}$ par enzyme par ribosome~; 30 par
ribosome par heure, 600 par ARNm par heure.
\textbf{16.} $100\times 600 = 60\,000$ enzymes par heure~; $10^{5}$ au
bout d'environ \qty{1.7}{h}.
\textbf{17.} Dix à vingt minutes pour que le cortisol entre, se fixe et
atteigne ses gènes, une heure avant que les premiers ARNm soient
nombreux et traduits, puis deux heures d'accumulation~: environ quatre
heures.
\textbf{18.} Une cascade modifie l'activité d'enzymes qui existent~; la
tâche du cortisol est de changer les enzymes qui existent, pour des
jours, ce que seule la transcription peut faire. Une frayeur ne peut pas
attendre quatre heures de nouvelles enzymes.
\textbf{19.} Plus de récepteurs, c'est plus de récepteurs occupés à
toute concentration d'adrénaline~: le taux d'occupation est inchangé
mais la réponse à chaque concentration est plus grande --- la courbe est
dilatée vers le haut, non décalée. Une \omterm{def:b2:cell-signalling:kinds}{hormone} lente règle le gain d'une
\omterm{def:b2:cell-signalling:kinds}{hormone} rapide.
\textbf{20.} Le glucose entre dans la cellule $\beta$ et y est
métabolisé~; l'ATP augmente et ferme un canal potassique sensible à
l'ATP~; la membrane se dépolarise~; des canaux calciques
voltage-dépendants s'ouvrent~; le calcium déclenche l'exocytose des
granules d'insuline --- le métabolisme couplé à la sécrétion par une
étape électrique.
\textbf{21.} Les deux récepteurs déclenchent deux voies qui convergent
sur les mêmes enzymes~: la kinase A du glucagon les phosphoryle, la
cascade de l'insuline active la phosphatase qui les déphosphoryle~;
l'état des enzymes, et donc le sens du métabolisme du glycogène, suit le
rapport des deux.
\textbf{22.} Glycémie à jeun élevée (le foie, sans frein, continue d'en
fabriquer), glycémie après un repas très élevée (pas d'absorption par
les muscles et le tissu adipeux)~; sans insuline, les réserves de
graisse libèrent des acides gras, que le foie convertit en corps
cétoniques acides --- l'acidocétose.
\textbf{23.} Un muscle au travail absorbe du glucose sans insuline~: au
cours d'une longue course, la glycémie baisse tandis que l'insuline
injectée, non régulée, continue de freiner le foie~: hypoglycémie et
malaise.
\textbf{24.} Boucle glycémique~: capteurs, les cellules $\beta$ et
$\alpha$~; deux \omterm{def:b2:cell-signalling:kinds}{hormones} effectrices pour les deux sens~; un gain élevé
qui maintient le niveau à un cinquième près~; des minutes pour agir.
\omterm{def:b2:blood-pressure:baroreflex}{Baroréflexe}~: des récepteurs à l'étirement, le \omterm{def:b2:heart:anatomy}{cœur} et les vaisseaux
comme effecteurs, un gain d'environ quatre, une seconde pour agir.
\textbf{25.} Taux d'occupation de $0.17$ après la libération~; gain de
$5\times 10^{9}$ par seconde par récepteur occupé~; production du foie
\qty{5}{g/min} en réalité~; le cortisol met environ quatre heures.
\end{solution}
